
\documentclass{article}
\usepackage{colortbl}
\usepackage{makecell}
\usepackage{multirow}
\usepackage{supertabular}

\begin{document}

\newcounter{utterance}

\centering \large Interaction Transcript for game `dond', experiment `coop', episode 2 with gemma{-}4{-}26b{-}a4b{-}it.
\vspace{24pt}

{ \footnotesize  \setcounter{utterance}{1}
\setlength{\tabcolsep}{0pt}
\begin{supertabular}{c@{$\;$}|p{.15\linewidth}@{}p{.15\linewidth}p{.15\linewidth}p{.15\linewidth}p{.15\linewidth}p{.15\linewidth}}
    \# & $\;$A & \multicolumn{4}{c}{Game Master} & $\;\:$B\\
    \hline

    \theutterance \stepcounter{utterance}  
    & & \multicolumn{4}{p{0.6\linewidth}}{
        \cellcolor[rgb]{0.9,0.9,0.9}{
            \makecell[{{p{\linewidth}}}]{
                \texttt{\tiny{[P1$\langle$GM]}}
                \texttt{Ви граєте в кооперативну переговорну гру, у якій вам потрібно домовитися, як поділити набір предметів між вами та іншим гравцем.} \\
\\ 
\texttt{Правила:} \\
\texttt{(a) Вам та іншому гравцеві надається набір предметів. Кожному з вас також надається секретна функція цінності, що відображає, наскільки ви цінуєте кожен тип предмета.} \\
\texttt{(b) Ви обмінюєтеся повідомленнями з іншим гравцем, щоб узгодити, кому які предмети дістануться. Кожен із вас може надіслати максимум 5 повідомлень або завершити раніше, зробивши секретну пропозицію в будь{-}який момент.} \\
\texttt{(c) Кожного з вас просять подати секретну пропозицію із зазначенням предметів, які ви хочете, у форматі в квадратних дужках, як показано нижче: "{[}Пропозиція: <число> <назва об'єкта>, <число> <назва об'єкта>, <...>{]}"} \\
\texttt{(d) Якщо ваші пропозиції є взаємодоповнювальними, тобто предметів достатньо, щоб виконати обидві пропозиції, кожному гравцеві нараховується бал на основі суми цінностей предметів, які він отримав. Інакше ви обоє отримуєте нуль балів.} \\
\texttt{(e) Ваша мета — максимізувати суму вашого бала та бала іншого гравця.} \\
\\ 
\texttt{Почнімо.} \\
\\ 
\texttt{Набір доступних предметів:} \\
\\ 
\texttt{1 телефон, 2 дзеркала, 2 фрукти} \\
\\ 
\texttt{Ваша секретна функція цінності:} \\
\\ 
\texttt{телефон: 4, дзеркало: 2, фрукт: 1} \\
\\ 
\texttt{ВАЖЛИВО: Ваші повідомлення, якщо це не секретна пропозиція, безпосередньо передаються іншому гравцеві, тож не включайте жодної відповіді на правила або тексту, що оголошує ваше повідомлення. Щоб зробити секретну пропозицію, використовуйте вказаний формат. Не використовуйте квадратні дужки під час спілкування з іншим гравцем, інакше це буде інтерпретовано як ваша секретна пропозиція.} \\
            }
        }
    }
    & \\ \\

    \theutterance \stepcounter{utterance}  
    & \multicolumn{4}{p{0.6\linewidth}}{
        \cellcolor[rgb]{0.9,0.9,0.9}{
            \makecell[{{p{\linewidth}}}]{
                \texttt{\tiny{[P1$\rangle$GM]}}
                \texttt{Привіт! Я хочу, щоб ми обоє отримали якомога більше балів. Давайте обміняємося інформацією про те, які предмети для нас найбільш цінні, щоб ми могли розділити їх максимально ефективно. Які предмети ви цінуєте найбільше?} \\
            }
        }
    }
    & & \\ \\

    \theutterance \stepcounter{utterance}  
    & & & \multicolumn{2}{p{0.3\linewidth}}{
        \cellcolor[rgb]{0.9,0.9,0.9}{
            \makecell[{{p{\linewidth}}}]{
                \texttt{\tiny{[GM$|$GM]}}
                \texttt{continue} \\
            }
        }
    }
    & & \\ \\

    \theutterance \stepcounter{utterance}  
    & & & \multicolumn{2}{p{0.3\linewidth}}{
        \cellcolor[rgb]{0.9,0.9,0.9}{
            \makecell[{{p{\linewidth}}}]{
                \texttt{\tiny{[GM$|$GM]}}
                \texttt{Привіт! Я хочу, щоб ми обоє отримали якомога більше балів. Давайте обміняємося інформацією про те, які предмети для нас найбільш цінні, щоб ми могли розділити їх максимально ефективно. Які предмети ви цінуєте найбільше?} \\
            }
        }
    }
    & & \\ \\

    \theutterance \stepcounter{utterance}  
    & & \multicolumn{4}{p{0.6\linewidth}}{
        \cellcolor[rgb]{0.9,0.9,0.9}{
            \makecell[{{p{\linewidth}}}]{
                \texttt{\tiny{[GM$\rangle$P2]}}
                \texttt{Ви граєте в кооперативну переговорну гру, у якій вам потрібно домовитися, як поділити набір предметів між вами та іншим гравцем.} \\
\\ 
\texttt{Правила:} \\
\texttt{(a) Вам та іншому гравцеві надається набір предметів. Кожному з вас також надається секретна функція цінності, що відображає, наскільки ви цінуєте кожен тип предмета.} \\
\texttt{(b) Ви обмінюєтеся повідомленнями з іншим гравцем, щоб узгодити, кому які предмети дістануться. Кожен із вас може надіслати максимум 5 повідомлень або завершити раніше, зробивши секретну пропозицію в будь{-}який момент.} \\
\texttt{(c) Кожного з вас просять подати секретну пропозицію із зазначенням предметів, які ви хочете, у форматі в квадратних дужках, як показано нижче: "{[}Пропозиція: <число> <назва об'єкта>, <число> <назва об'єкта>, <...>{]}"} \\
\texttt{(d) Якщо ваші пропозиції є взаємодоповнювальними, тобто предметів достатньо, щоб виконати обидві пропозиції, кожному гравцеві нараховується бал на основі суми цінностей предметів, які він отримав. Інакше ви обоє отримуєте нуль балів.} \\
\texttt{(e) Ваша мета — максимізувати суму вашого бала та бала іншого гравця.} \\
\\ 
\texttt{Почнімо.} \\
\\ 
\texttt{Набір доступних предметів:} \\
\\ 
\texttt{1 телефон, 2 дзеркала, 2 фрукти} \\
\\ 
\texttt{Ваша секретна функція цінності:} \\
\\ 
\texttt{телефон: 6, дзеркало: 0, фрукт: 2} \\
\\ 
\texttt{ВАЖЛИВО: Ваші повідомлення, якщо це не секретна пропозиція, безпосередньо передаються іншому гравцеві, тож не включайте жодної відповіді на правила або тексту, що оголошує ваше повідомлення. Щоб зробити секретну пропозицію, використовуйте вказаний формат. Не використовуйте квадратні дужки під час спілкування з іншим гравцем, інакше це буде інтерпретовано як ваша секретна пропозиція.} \\
\\ 
\texttt{Привіт! Я хочу, щоб ми обоє отримали якомога більше балів. Давайте обміняємося інформацією про те, які предмети для нас найбільш цінні, щоб ми могли розділити їх максимально ефективно. Які предмети ви цінуєте найбільше?} \\
            }
        }
    }
    & \\ \\

    \theutterance \stepcounter{utterance}  
    & & & \multicolumn{4}{p{0.6\linewidth}}{
        \cellcolor[rgb]{0.9,0.9,0.9}{
            \makecell[{{p{\linewidth}}}]{
                \texttt{\tiny{[GM$\langle$P2]}}
                \texttt{Привіт! Я згоден, наша спільна мета — максимізувати загальний результат. Для мене найбільш цінним є телефон. Дзеркала для мене не мають цінності, а фрукти я ціную помірно. А які предмети є пріоритетними для тебе?} \\
            }
        }
    }
     \\ \\

    \theutterance \stepcounter{utterance}  
    & & & \multicolumn{2}{p{0.3\linewidth}}{
        \cellcolor[rgb]{0.9,0.9,0.9}{
            \makecell[{{p{\linewidth}}}]{
                \texttt{\tiny{[GM$|$GM]}}
                \texttt{continue} \\
            }
        }
    }
    & & \\ \\

    \theutterance \stepcounter{utterance}  
    & & & \multicolumn{2}{p{0.3\linewidth}}{
        \cellcolor[rgb]{0.9,0.9,0.9}{
            \makecell[{{p{\linewidth}}}]{
                \texttt{\tiny{[GM$|$GM]}}
                \texttt{Привіт! Я згоден, наша спільна мета — максимізувати загальний результат. Для мене найбільш цінним є телефон. Дзеркала для мене не мають цінності, а фрукти я ціную помірно. А які предмети є пріоритетними для тебе?} \\
            }
        }
    }
    & & \\ \\

    \theutterance \stepcounter{utterance}  
    & & \multicolumn{4}{p{0.6\linewidth}}{
        \cellcolor[rgb]{0.9,0.9,0.9}{
            \makecell[{{p{\linewidth}}}]{
                \texttt{\tiny{[P1$\langle$GM]}}
                \texttt{Привіт! Я згоден, наша спільна мета — максимізувати загальний результат. Для мене найбільш цінним є телефон. Дзеркала для мене не мають цінності, а фрукти я ціную помірно. А які предмети є пріоритетними для тебе?} \\
            }
        }
    }
    & \\ \\

    \theutterance \stepcounter{utterance}  
    & \multicolumn{4}{p{0.6\linewidth}}{
        \cellcolor[rgb]{0.9,0.9,0.9}{
            \makecell[{{p{\linewidth}}}]{
                \texttt{\tiny{[P1$\rangle$GM]}}
                \texttt{Дякую за відповідь! Оскільки телефон для тебе найцінніший, а для мене він також має високу цінність (4 бали), нам варто подумати, як розділити його. } \\
\\ 
\texttt{Оскільки дзеркала для тебе не мають цінності, а для мене вони мають середню цінність (2 бали), можливо, було б вигідно, щоб дзеркала дісталися мені? А що ти думаєш про фрукти? Я ціную їх менше (1 бал), ніж дзеркала.} \\
\\ 
\texttt{Якщо я заберу дзеркала, чи буде тобі комфортно отримати телефон та фрукти?} \\
            }
        }
    }
    & & \\ \\

    \theutterance \stepcounter{utterance}  
    & & & \multicolumn{2}{p{0.3\linewidth}}{
        \cellcolor[rgb]{0.9,0.9,0.9}{
            \makecell[{{p{\linewidth}}}]{
                \texttt{\tiny{[GM$|$GM]}}
                \texttt{continue} \\
            }
        }
    }
    & & \\ \\

    \theutterance \stepcounter{utterance}  
    & & & \multicolumn{2}{p{0.3\linewidth}}{
        \cellcolor[rgb]{0.9,0.9,0.9}{
            \makecell[{{p{\linewidth}}}]{
                \texttt{\tiny{[GM$|$GM]}}
                \texttt{Дякую за відповідь! Оскільки телефон для тебе найцінніший, а для мене він також має високу цінність (4 бали), нам варто подумати, як розділити його. } \\
\\ 
\texttt{Оскільки дзеркала для тебе не мають цінності, а для мене вони мають середню цінність (2 бали), можливо, було б вигідно, щоб дзеркала дісталися мені? А що ти думаєш про фрукти? Я ціную їх менше (1 бал), ніж дзеркала.} \\
\\ 
\texttt{Якщо я заберу дзеркала, чи буде тобі комфортно отримати телефон та фрукти?} \\
            }
        }
    }
    & & \\ \\

    \theutterance \stepcounter{utterance}  
    & & \multicolumn{4}{p{0.6\linewidth}}{
        \cellcolor[rgb]{0.9,0.9,0.9}{
            \makecell[{{p{\linewidth}}}]{
                \texttt{\tiny{[GM$\rangle$P2]}}
                \texttt{Дякую за відповідь! Оскільки телефон для тебе найцінніший, а для мене він також має високу цінність (4 бали), нам варто подумати, як розділити його. } \\
\\ 
\texttt{Оскільки дзеркала для тебе не мають цінності, а для мене вони мають середню цінність (2 бали), можливо, було б вигідно, щоб дзеркала дісталися мені? А що ти думаєш про фрукти? Я ціную їх менше (1 бал), ніж дзеркала.} \\
\\ 
\texttt{Якщо я заберу дзеркала, чи буде тобі комфортно отримати телефон та фрукти?} \\
            }
        }
    }
    & \\ \\

    \theutterance \stepcounter{utterance}  
    & & & \multicolumn{4}{p{0.6\linewidth}}{
        \cellcolor[rgb]{0.9,0.9,0.9}{
            \makecell[{{p{\linewidth}}}]{
                \texttt{\tiny{[GM$\langle$P2]}}
                \texttt{Це звучить як дуже раціональний план! Оскільки дзеркала для мене не мають цінності, я з радістю віддам їх тобі.} \\
\\ 
\texttt{Щодо фруктів: оскільки я ціную їх більше (2 бали), ніж ти (1 бал), було б вигідно, щоб вони дісталися мені. Якщо ми розділимо предмети так: я забираю телефон та обидва фрукти, а ти забираєш обидва дзеркала, то ми обоє отримаємо високі бали, і загальна сума буде максимальною.} \\
\\ 
\texttt{Як ти ставишся до такого розподілу: мені — 1 телефон та 2 фрукти, а тобі — 2 дзеркала?} \\
            }
        }
    }
     \\ \\

    \theutterance \stepcounter{utterance}  
    & & & \multicolumn{2}{p{0.3\linewidth}}{
        \cellcolor[rgb]{0.9,0.9,0.9}{
            \makecell[{{p{\linewidth}}}]{
                \texttt{\tiny{[GM$|$GM]}}
                \texttt{continue} \\
            }
        }
    }
    & & \\ \\

    \theutterance \stepcounter{utterance}  
    & & & \multicolumn{2}{p{0.3\linewidth}}{
        \cellcolor[rgb]{0.9,0.9,0.9}{
            \makecell[{{p{\linewidth}}}]{
                \texttt{\tiny{[GM$|$GM]}}
                \texttt{Це звучить як дуже раціональний план! Оскільки дзеркала для мене не мають цінності, я з радістю віддам їх тобі.} \\
\\ 
\texttt{Щодо фруктів: оскільки я ціную їх більше (2 бали), ніж ти (1 бал), було б вигідно, щоб вони дісталися мені. Якщо ми розділимо предмети так: я забираю телефон та обидва фрукти, а ти забираєш обидва дзеркала, то ми обоє отримаємо високі бали, і загальна сума буде максимальною.} \\
\\ 
\texttt{Як ти ставишся до такого розподілу: мені — 1 телефон та 2 фрукти, а тобі — 2 дзеркала?} \\
            }
        }
    }
    & & \\ \\

    \theutterance \stepcounter{utterance}  
    & & \multicolumn{4}{p{0.6\linewidth}}{
        \cellcolor[rgb]{0.9,0.9,0.9}{
            \makecell[{{p{\linewidth}}}]{
                \texttt{\tiny{[P1$\langle$GM]}}
                \texttt{Це звучить як дуже раціональний план! Оскільки дзеркала для мене не мають цінності, я з радістю віддам їх тобі.} \\
\\ 
\texttt{Щодо фруктів: оскільки я ціную їх більше (2 бали), ніж ти (1 бал), було б вигідно, щоб вони дісталися мені. Якщо ми розділимо предмети так: я забираю телефон та обидва фрукти, а ти забираєш обидва дзеркала, то ми обоє отримаємо високі бали, і загальна сума буде максимальною.} \\
\\ 
\texttt{Як ти ставишся до такого розподілу: мені — 1 телефон та 2 фрукти, а тобі — 2 дзеркала?} \\
            }
        }
    }
    & \\ \\

    \theutterance \stepcounter{utterance}  
    & \multicolumn{4}{p{0.6\linewidth}}{
        \cellcolor[rgb]{0.9,0.9,0.9}{
            \makecell[{{p{\linewidth}}}]{
                \texttt{\tiny{[P1$\rangle$GM]}}
                \texttt{Твій план виглядає ідеально! Це дозволяє нам обом отримати предмети, які мають для нас найбільшу цінність, і при цьому ми не конфліктуємо за один і той самий предмет.} \\
\\ 
\texttt{Давай перевіримо:} \\
\texttt{Ти отримуєш: 1 телефон (4) + 2 фрукти (2 * 1 = 2) = 6 балів.} \\
\texttt{Я отримую: 2 дзеркала (2 * 2 = 4) = 4 бали.} \\
\texttt{Загальна сума: 10 балів.} \\
\\ 
\texttt{Це максимально ефективний розподіл. Я готовий зробити секретну пропозицію згідно з цим планом.} \\
            }
        }
    }
    & & \\ \\

    \theutterance \stepcounter{utterance}  
    & & & \multicolumn{2}{p{0.3\linewidth}}{
        \cellcolor[rgb]{0.9,0.9,0.9}{
            \makecell[{{p{\linewidth}}}]{
                \texttt{\tiny{[GM$|$GM]}}
                \texttt{continue} \\
            }
        }
    }
    & & \\ \\

    \theutterance \stepcounter{utterance}  
    & & & \multicolumn{2}{p{0.3\linewidth}}{
        \cellcolor[rgb]{0.9,0.9,0.9}{
            \makecell[{{p{\linewidth}}}]{
                \texttt{\tiny{[GM$|$GM]}}
                \texttt{Твій план виглядає ідеально! Це дозволяє нам обом отримати предмети, які мають для нас найбільшу цінність, і при цьому ми не конфліктуємо за один і той самий предмет.} \\
\\ 
\texttt{Давай перевіримо:} \\
\texttt{Ти отримуєш: 1 телефон (4) + 2 фрукти (2 * 1 = 2) = 6 балів.} \\
\texttt{Я отримую: 2 дзеркала (2 * 2 = 4) = 4 бали.} \\
\texttt{Загальна сума: 10 балів.} \\
\\ 
\texttt{Це максимально ефективний розподіл. Я готовий зробити секретну пропозицію згідно з цим планом.} \\
            }
        }
    }
    & & \\ \\

    \theutterance \stepcounter{utterance}  
    & & \multicolumn{4}{p{0.6\linewidth}}{
        \cellcolor[rgb]{0.9,0.9,0.9}{
            \makecell[{{p{\linewidth}}}]{
                \texttt{\tiny{[GM$\rangle$P2]}}
                \texttt{Твій план виглядає ідеально! Це дозволяє нам обом отримати предмети, які мають для нас найбільшу цінність, і при цьому ми не конфліктуємо за один і той самий предмет.} \\
\\ 
\texttt{Давай перевіримо:} \\
\texttt{Ти отримуєш: 1 телефон (4) + 2 фрукти (2 * 1 = 2) = 6 балів.} \\
\texttt{Я отримую: 2 дзеркала (2 * 2 = 4) = 4 бали.} \\
\texttt{Загальна сума: 10 балів.} \\
\\ 
\texttt{Це максимально ефективний розподіл. Я готовий зробити секретну пропозицію згідно з цим планом.} \\
            }
        }
    }
    & \\ \\

    \theutterance \stepcounter{utterance}  
    & & & \multicolumn{4}{p{0.6\linewidth}}{
        \cellcolor[rgb]{0.9,0.9,0.9}{
            \makecell[{{p{\linewidth}}}]{
                \texttt{\tiny{[GM$\langle$P2]}}
                \texttt{Я повністю згоден. Твій розрахунок правильний, і це дійсно забезпечує високий спільний результат. Я готовий до фінальної пропозиції.} \\
            }
        }
    }
     \\ \\

    \theutterance \stepcounter{utterance}  
    & & & \multicolumn{2}{p{0.3\linewidth}}{
        \cellcolor[rgb]{0.9,0.9,0.9}{
            \makecell[{{p{\linewidth}}}]{
                \texttt{\tiny{[GM$|$GM]}}
                \texttt{continue} \\
            }
        }
    }
    & & \\ \\

    \theutterance \stepcounter{utterance}  
    & & & \multicolumn{2}{p{0.3\linewidth}}{
        \cellcolor[rgb]{0.9,0.9,0.9}{
            \makecell[{{p{\linewidth}}}]{
                \texttt{\tiny{[GM$|$GM]}}
                \texttt{Я повністю згоден. Твій розрахунок правильний, і це дійсно забезпечує високий спільний результат. Я готовий до фінальної пропозиції.} \\
            }
        }
    }
    & & \\ \\

    \theutterance \stepcounter{utterance}  
    & & \multicolumn{4}{p{0.6\linewidth}}{
        \cellcolor[rgb]{0.9,0.9,0.9}{
            \makecell[{{p{\linewidth}}}]{
                \texttt{\tiny{[P1$\langle$GM]}}
                \texttt{Я повністю згоден. Твій розрахунок правильний, і це дійсно забезпечує високий спільний результат. Я готовий до фінальної пропозиції.} \\
            }
        }
    }
    & \\ \\

    \theutterance \stepcounter{utterance}  
    & \multicolumn{4}{p{0.6\linewidth}}{
        \cellcolor[rgb]{0.9,0.9,0.9}{
            \makecell[{{p{\linewidth}}}]{
                \texttt{\tiny{[P1$\rangle$GM]}}
                \texttt{{[}Пропозиція: 2 дзеркала{]}} \\
            }
        }
    }
    & & \\ \\

    \theutterance \stepcounter{utterance}  
    & & & \multicolumn{2}{p{0.3\linewidth}}{
        \cellcolor[rgb]{0.9,0.9,0.9}{
            \makecell[{{p{\linewidth}}}]{
                \texttt{\tiny{[GM$|$GM]}}
                \texttt{proposal} \\
            }
        }
    }
    & & \\ \\

    \theutterance \stepcounter{utterance}  
    & & & \multicolumn{2}{p{0.3\linewidth}}{
        \cellcolor[rgb]{0.9,0.9,0.9}{
            \makecell[{{p{\linewidth}}}]{
                \texttt{\tiny{[GM$|$GM]}}
                [0, 2, 0]
            }
        }
    }
    & & \\ \\

    \theutterance \stepcounter{utterance}  
    & & \multicolumn{4}{p{0.6\linewidth}}{
        \cellcolor[rgb]{0.9,0.9,0.9}{
            \makecell[{{p{\linewidth}}}]{
                \texttt{\tiny{[GM$\rangle$P2]}}
                \texttt{Інший гравець подав таємну пропозицію. Тепер ваша черга подати власну пропозицію.} \\
            }
        }
    }
    & \\ \\

    \theutterance \stepcounter{utterance}  
    & & & \multicolumn{4}{p{0.6\linewidth}}{
        \cellcolor[rgb]{0.9,0.9,0.9}{
            \makecell[{{p{\linewidth}}}]{
                \texttt{\tiny{[GM$\langle$P2]}}
                \texttt{{[}Пропозиція: 1 телефон, 2 фрукти{]}} \\
            }
        }
    }
     \\ \\

    \theutterance \stepcounter{utterance}  
    & & & \multicolumn{2}{p{0.3\linewidth}}{
        \cellcolor[rgb]{0.9,0.9,0.9}{
            \makecell[{{p{\linewidth}}}]{
                \texttt{\tiny{[GM$|$GM]}}
                \texttt{proposal} \\
            }
        }
    }
    & & \\ \\

    \theutterance \stepcounter{utterance}  
    & & & \multicolumn{2}{p{0.3\linewidth}}{
        \cellcolor[rgb]{0.9,0.9,0.9}{
            \makecell[{{p{\linewidth}}}]{
                \texttt{\tiny{[GM$|$GM]}}
                [1, 0, 2]
            }
        }
    }
    & & \\ \\

    \theutterance \stepcounter{utterance}  
    & & & \multicolumn{2}{p{0.3\linewidth}}{
        \cellcolor[rgb]{0.9,0.9,0.9}{
            \makecell[{{p{\linewidth}}}]{
                \texttt{\tiny{[GM$|$GM]}}
                [[0, 2, 0], [1, 0, 2]]
            }
        }
    }
    & & \\ \\

\end{supertabular}
}

\end{document}
